\subsection{Purpose}

This section defines requirements for securely retiring, replacing, transferring, or removing systems, applications, accounts, equipment, services, and information.

The purpose is to prevent unauthorized access, data exposure, service disruption, loss of records, and unmanaged dependencies when a system or service is no longer required.

\subsection{Scope}

These requirements apply to:
\begin{itemize}
    \item Servers, workstations, laptops, mobile devices, and network equipment
    \item Virtual machines, containers, cloud resources, and storage volumes
    \item Applications, databases, websites, and hosted services
    \item User, administrator, service, vendor, and emergency accounts
    \item Encryption keys, certificates, API keys, tokens, and credentials
    \item Backups, logs, documentation, and other associated records
\end{itemize}

\subsection{Decommissioning Authorization}

A system or service \must{} have an approved decommissioning plan before it is retired or removed.

The plan \must{} identify:
\begin{itemize}
    \item The system, service, device, or account being decommissioned
    \item The business and technical owner
    \item The reason for decommissioning
    \item The planned date and responsible personnel
    \item Dependencies and connected systems
    \item Information and records that must be retained
    \item Backup and recovery considerations
    \item Data sanitization and disposal methods
    \item Required notifications and approvals
\end{itemize}

Critical systems \must{} not be decommissioned without approval from the system owner and the appropriate management or technical authority.

\subsection{Dependency Review}

Before decommissioning, the responsible administrator \must{} identify systems, applications, integrations, users, vendors, network connections, scheduled tasks, certificates, credentials, backups, and other dependencies.

The dependency review \should{} include:
\begin{itemize}
    \item Applications and services that rely on the system
    \item Network routes, firewall rules, and integrations
    \item Data feeds and automated processes
    \item Monitoring and alerting configurations
    \item Backup and replication jobs
    \item DNS records and certificates
    \item Service accounts and stored credentials
    \item Documentation and support procedures
\end{itemize}

Dependencies \must{} be migrated, disabled, or formally accepted before final removal.

\subsection{Data Retention and Transfer}

Information stored on a system \must{} be reviewed before decommissioning.

The system owner \must{} determine whether information should be:
\begin{itemize}
    \item Transferred to a replacement system
    \item Archived according to retention requirements
    \item Included in an approved backup
    \item Securely deleted
    \item Returned to a customer, vendor, or other authorized party
\end{itemize}

Data \must{} not be retained longer than required unless a documented business, legal, contractual, or regulatory reason exists.

Archived information \must{} remain subject to access control, classification, encryption, retention, monitoring, and disposal requirements.

\subsection{Account and Credential Removal}

Accounts associated with a decommissioned system \must{} be disabled or removed when they are no longer required.

This includes:
\begin{itemize}
    \item User accounts
    \item Administrator accounts
    \item Service accounts
    \item Vendor accounts
    \item Emergency accounts
    \item API keys and access tokens
    \item SSH keys and certificates
    \item Stored passwords and connection credentials
\end{itemize}

Credentials that may have been exposed during the system's lifecycle \should{} be rotated or revoked.

Network access rules, firewall exceptions, VPN access, remote-management settings, and other access paths \must{} be removed when no longer required.

\subsection{Data Sanitization}

Storage media \must{} be sanitized before equipment is reused, transferred, returned, sold, recycled, or disposed of.

Sanitization methods \must{} be appropriate to the sensitivity of the information and the type of media. Acceptable methods may include:
\begin{itemize}
    \item Approved secure erasure
    \item Cryptographic erasure
    \item Factory reset combined with secure deletion where appropriate
    \item Physical destruction
    \item Destruction by an approved disposal provider
\end{itemize}

Ordinary file deletion or reformatting \mustnot{} be considered sufficient for Confidential or Restricted information unless the approved sanitization process verifies that the information cannot reasonably be recovered.

\subsection{Hardware and Media Disposal}

Hardware and removable media \must{} be tracked through disposal or transfer.

The disposal record \should{} include:
\begin{itemize}
    \item Asset identifier
    \item Equipment description
    \item Assigned owner or department
    \item Date of disposal or transfer
    \item Sanitization method
    \item Person or provider performing the work
    \item Final disposition
\end{itemize}

Third-party disposal providers \must{} be approved and \should{} provide evidence of secure disposal when required by the data classification or contract.

\subsection{Service and Cloud Decommissioning}

Cloud resources and hosted services \must{} be reviewed before cancellation or deletion.

The review \must{} address:
\begin{itemize}
    \item Data stored by the provider
    \item Snapshots, replicas, and backup copies
    \item Connected accounts and integrations
    \item Domains, certificates, and DNS records
    \item Billing and subscription ownership
    \item Audit logs and records required for retention
    \item Provider-specific deletion and recovery procedures
\end{itemize}

Cloud resources \must{} be deleted or disabled after required data, records, and dependencies have been handled.

\subsection{Verification}

After decommissioning, the responsible administrator \must{} verify that:
\begin{itemize}
    \item The system is no longer accessible through approved or known access paths
    \item Required data has been transferred or archived
    \item Unneeded data has been securely removed
    \item Accounts, credentials, keys, and certificates have been disabled or revoked
    \item Monitoring and backup jobs have been updated
    \item Network rules and integrations have been removed
    \item Asset and system inventories are current
    \item Documentation reflects the final state
\end{itemize}

Critical system decommissioning \should{} include a second-person review or formal approval of the completion evidence.

\subsection{Emergency Decommissioning}

A system may be isolated or decommissioned without following the normal schedule when necessary to address an active compromise, unacceptable risk, safety issue, legal requirement, or serious operational threat.

Emergency actions \must{} be documented after the immediate risk has been controlled.

The documentation \must{} include the reason for the emergency action, approvals or notifications, data-handling decisions, and follow-up activities.

\subsection{Records}

Decommissioning records \must{} be retained according to applicable business, legal, contractual, and records-management requirements.

The organization \should{} retain evidence of approvals, data transfers, sanitization, disposal, verification, and final inventory updates for critical systems and sensitive information.

\subsection{Review}

Decommissioning procedures \should{} be reviewed at least annually and whenever there is a significant change to technology, service providers, retention requirements, or organizational processes.